\chapter{Fundamentos teóricos}\label{ch:theoretical_background}

% ============================================================
%  BLOQUE A — CONTEXTO DINÁMICO: ¿QUÉ SON LOS CICLONES
%             EXTRATROPICALES EN EL ATLÁNTICO SUR?
%  Propósito: Establecer las características físicas generales
%  de los ciclones extratropicales en el Atlántico Sur,
%  incluyendo su diversidad estructural y la influencia
%  orográfica y oceánica que los distingue de la teoría
%  clásica. Es el punto de partida conceptual antes de
%  introducir las herramientas de detección (Bloque B).
% ============================================================

\section{Ciclones extratropicales en el Atlántico Sur}\label{sec:tb_fundamentos}

% --- A1: Teoría clásica y sus limitaciones regionales ---
La teoría clásica (e.g., modelo noruego) define a los ciclones extratropicales como vórtices de núcleo frío cuya génesis está dominada por la inestabilidad baroclínica de latitudes medias, perspectiva consolidada globalmente por las climatologías de \citet{hoskins2005new}. Sin embargo, la caracterización de estos sistemas en el Atlántico Sur requiere matices dinámicos adicionales. \citet{sinclair1995climatology} describió las características del ciclo de vida de los ciclones en el Hemisferio Sur, resaltando que la ciclogénesis está influenciada por los gradientes térmicos oceánicos y la orografía. De manera específica, \citet{inatsu2004zonal} demostraron que la asimetría zonal del \textit{storm track} en el Hemisferio Sur está forzada por la topografía de los Andes, la cual modula la propagación de ondas baroclínicas sobre el subcontinente sudamericano \citep{vera2002cold}, definiendo regiones preferenciales de desarrollo en latitudes subtropicales y extratropicales.

% --- A2: Diversidad estructural y procesos diabáticos ---
Esta variabilidad estructural implica que la definición adiabática estándar es insuficiente para la región. Si bien la inestabilidad baroclínica proporciona el entorno favorable para la génesis, la evolución y la intensidad del sistema puede ser modulada por procesos diabáticos asociados a la liberación de calor latente, como concluye \citet{coutodesouza2024thesis} en su revisión de la física de estos sistemas para el Atlántico Sur. Este comportamiento es ilustrado por los experimentos de sensibilidad numérica presentados por \citet{reboita2022from}, donde la supresión de los flujos de calor latente altera la estructura térmica y la intensidad del vórtice, evidenciando que la dinámica de estos sistemas es intrínsecamente híbrida.

% ============================================================
%  BLOQUE B — DETECCIÓN Y SEGUIMIENTO: VORTICIDAD RELATIVA
%             Y MARCO LAGRANGIANO
%  Propósito: Fundamentar las dos decisiones instrumentales
%  centrales de la metodología: (B1) el uso de la vorticidad
%  relativa ζ₈₅₀ como variable de tracking, y (B2) la
%  adopción del marco de referencia lagrangiano centrado en
%  el sistema. Ambas decisiones son prerequisito conceptual
%  para todo el diseño experimental posterior.
% ============================================================

\section{Detección lagrangiana de ciclones}\label{sec:tb_vorticidad}

% --- B1: Limitaciones de la presión y justificación de ζ₈₅₀ ---
Debido a la extensión oceánica del Hemisferio Sur, los ciclones extratropicales se desarrollan inmersos en un flujo zonal de los oestes. Esta corriente de fondo hace que los sistemas se trasladen con gran rapidez y que, en sus etapas iniciales, la caída de presión quede enmascarada, presentándose a menudo como vaguadas abiertas en lugar de centros cerrados \citep{sinclair1994objective}. Por esta razón, la presión al nivel del mar (MSLP) resulta una métrica inadecuada para la detección temprana.

Para resolver este problema, \citet{sinclair1997objective} recomienda utilizar la vorticidad relativa en 850 hPa ($\zeta_{850}$) como variable de seguimiento. Conceptualmente, la vorticidad relativa mide la rotación macroscópica local del fluido atmosférico respecto a un marco de referencia en la superficie terrestre. Físicamente se define como:

\begin{equation}\label{eq:vorticidad_relativa}
\zeta = \mathbf{k} \cdot (\nabla \times \mathbf{V}) = \frac{\partial v}{\partial x} - \frac{\partial u}{\partial y},
\end{equation}

\noindent donde $u$ y $v$ son las componentes zonal y meridional del viento horizontal, respectivamente. Según \citet{hoskins2005new}, el uso de $\zeta$ permite filtrar el flujo de gran escala y capturar la rotación inducida por forzantes de mesoescala antes de que se manifiesten en el campo de presión. Esta capacidad es validada por \citet{gramcianinov2019properties}, quienes utilizan $\zeta_{850}$ para identificar con precisión las regiones de ciclogénesis en el Atlántico Sur, y por \citet{padilhareinke2024objective}, quienes demuestran que los algoritmos basados en este nivel isobárico son más eficientes para capturar la posición del ciclón durante su etapa de mayor intensidad.

Recientemente, el uso de $\zeta_{850}$ se ha extendido hacia la caracterización estructural: \citet{coutodesouza2024new} demuestran que la evolución temporal de esta variable es el predictor más confiable para segmentar objetivamente el ciclo de vida en fases dinámicas. En un contexto global, \citet{corner2025classification} ratifican a la vorticidad relativa en 850 hPa como una de las métricas irreductibles para describir la intensidad y la estructura de los ciclones extratropicales.

% --- B2: Estratificación por intensidad mediante percentiles de ζ₈₅₀ ---
El uso de umbrales percentílicos sobre la distribución de $\zeta_{850}$ para estratificar ciclones por intensidad dinámica cuenta con respaldo metodológico sólido. \citet{priestley2022improved} aplican explícitamente este criterio, seleccionando el 10\% de ciclones con mayor vorticidad ciclónica en su pico de intensidad para construir compositos representativos de la estructura ciclónica. En esta tesis, ese subconjunto se denominará p90 (percentil 90 de vorticidad absoluta), ya que aísla los sistemas más intensos desde el punto de vista dinámico. De forma complementaria, \citet{andrade2024composite} documentan para el Atlántico Sudoccidental que los ciclones intensos seleccionados por percentil de presión mínima exhiben patrones sinópticos diferenciados y mayor coherencia entre los extremos de viento en superficie y la intensidad del vórtice. Aunque presión y vorticidad miden aspectos distintos de la intensidad ciclónica, ambos umbrales percentílicos responden a la misma lógica: retener solo los casos más intensos, lo que permite validar estas submuestras como poblaciones física y estructuralmente homogéneas.

% --- B3: Marco de referencia lagrangiano ---
La adopción de un marco de referencia lagrangiano centrado en el sistema —en contraposición a un análisis euleriano de campo fijo— es la decisión instrumental que hace posible comparar la estructura interna de ciclones que ocurren en distintos momentos y posiciones geográficas. En el enfoque euleriano, los campos de viento y precipitación de distintos sistemas se superponen en coordenadas geográficas absolutas, diluyendo las asimetrías estructurales propias de cada ciclón en el promedio geográfico. El enfoque lagrangiano, en cambio, centra cada sistema en su propio núcleo de vorticidad y opera en coordenadas relativas $(\Delta x, \Delta y)$, preservando la organización espacial interna del ciclón independientemente de su posición geográfica instantánea. Esta perspectiva, fundamentada en los principios de seguimiento objetivo de \citet{sinclair1994objective}, permite que la composición de múltiples eventos revele patrones estructurales coherentes que serían invisibles en un análisis de campo fijo. Es precisamente este marco el que hace posible definir un dominio móvil de análisis, centrar los campos horarios en cada paso de tiempo de la trayectoria, y construir las distribuciones espaciales de probabilidad de los máximos que constituyen el núcleo del análisis estadístico de esta tesis.

% ============================================================
%  BLOQUE C — REGIONES DE CICLOGÉNESIS
%  Propósito: Describir los tres focos de ciclogénesis
%  (SBR, LPB, ARG) y sus mecanismos físicos diferenciados.
%  Este bloque se ubica después de la vorticidad porque
%  es precisamente ζ₈₅₀ el instrumento que permite delimitar
%  estas regiones objetivamente (Reboita 2010, Gramcianinov
%  2019). Su posición aquí establece el contexto geográfico
%  antes de introducir los modelos conceptuales de estructura
%  interna (Bloque D).
% ============================================================

\section{Regiones de ciclogénesis}\label{sec:tb_regiones}

% --- C0: Delimitación general de las regiones ---
A partir del análisis de la vorticidad relativa en niveles bajos, \citet{reboita2010south} establecieron la zonificación climática de la ciclogénesis en el Atlántico Sur, delimitando la costa de Uruguay y sur de Brasil, y la región costera de Argentina, como los principales focos de actividad. La interacción entre el flujo de los oestes, la topografía de los Andes y la Temperatura Superficial del Mar (TSM) define tres regiones de génesis con características físicas diferenciadas \citep{gramcianinov2019properties}, cuya correspondencia espacial con las regiones SBR, LPB y ARG adoptadas en esta tesis fue establecida en la Sección~\ref{sec:definicion_regiones} de la Introducción.

\subsection{Sur-sureste de Brasil (SBR)}\label{subsec:tb_sbr}

% --- C1: Mecanismo híbrido de SBR ---
Según \citet{reboita2022from}, la región SBR constituye una de las regiones preferenciales de ciclogénesis en la costa sur/sudeste de Brasil. Los sistemas que actúan en este sector pueden presentar características subtropicales durante su ciclo de vida, siendo denominados ciclones ``híbridos'': aunque se inician por inestabilidad baroclínica, su intensificación está modulada por procesos diabáticos y la inestabilidad de la capa límite marina, lo que favorece el desarrollo de convección en el sector cálido \citep{marrafon2022classificacao}. El análisis energético de \citet{coutodesouza2024thesis} revela que, a diferencia de las regiones australes, los ciclones en SBR presentan una conversión de energía mixta, donde los procesos diabáticos (liberación de calor latente) juegan un rol tan importante como la baroclinicidad. \citet{gramcianinov2024early} demuestran además que, durante la ciclogénesis temprana en SBR, los modelos de alta resolución logran representar las estructuras de mesoescala de advección cálida y convergencia de humedad en capas bajas, validando que la dinámica híbrida de SBR no es un artefacto estadístico sino una señal física robusta.

\subsection{Cuenca del Río de la Plata (LPB)}\label{subsec:tb_lpb}

% --- C2: Ciclogénesis de sotavento y rol del SALLJ en LPB ---
La región comprendida entre $30^\circ$S y $40^\circ$S constituye uno de los principales \textit{hotspots} de ciclogénesis en el Atlántico Sudoccidental. Su mecanismo dominante es la ciclogénesis de sotavento: el flujo de los oestes experimenta una compresión y posterior estiramiento vertical de la columna de atmósfera al atravesar los Andes y, por conservación de la vorticidad potencial, este estiramiento en el flanco oriental induce una caída de presión en superficie que favorece la iniciación de la rotación ciclónica \citep{gan1994influence}. \citet{reboita2010south} muestran que esta región se caracteriza por una alta frecuencia de sistemas ciclónicos y una interacción con el transporte de humedad desde latitudes tropicales, mientras que el aporte de aire cálido y húmedo asociado al Chorro de Capas Bajas de Sudamérica (SALLJ) resulta determinante para la intensidad de la precipitación \citep{crespo2020potential}. \citet{gramcianinov2024early} confirman mediante modelos regionales que la convergencia de humedad en capas bajas constituye el mecanismo dominante de intensificación en LPB, proyectando además que bajo escenarios de calentamiento global los sistemas que persistan exhibirán flujos de humedad y advección de calor más intensos en sus etapas tempranas.

\subsection{Argentina (ARG)}\label{subsec:tb_arg}

% --- C3: Régimen baroclínico clásico de ARG ---
El sector austral ($>40^\circ$S) responde a un régimen dinámico clásico dominado por la baroclinicidad de gran escala. \citet{hoskins2005new} demuestran que el desarrollo en estas latitudes está estrictamente acoplado en toda la columna troposférica, siendo gobernado por la inestabilidad baroclínica asociada al chorro polar y la propagación de ondas de Rossby a lo largo del \textit{storm track} de alta latitud. \citet{simmonds2000variability} indican que los sistemas ciclónicos en ARG dependen menos de forzantes locales de superficie y más de la dinámica de la alta troposfera, lo que genera una evolución más rápida y directamente acoplada al flujo de niveles altos, en contraste con las regiones subtropicales del sudeste de Brasil \citep{vera2002cold}.

% ============================================================
%  BLOQUE D — ESTRUCTURA INTERNA: MODELOS CONCEPTUALES
%             Y CICLO DE VIDA
%  Propósito: Describir los modelos conceptuales que explican
%  la organización tridimensional del ciclón (Shapiro-Keyser,
%  WCB, CCB) y la segmentación objetiva de su ciclo de vida
%  en fases. Este bloque es el puente conceptual entre la
%  detección (Bloque B) y la distribución de los extremos
%  en superficie (Bloque E).
% ============================================================

\section{Ciclo de vida y modelos conceptuales}\label{sec:tb_evolucion}

\subsection{Modelo conceptual de estructura interna}\label{subsec:tb_modelos}

% --- D1: Modelo de Shapiro-Keyser como paradigma regional ---
La distribución espacial de las variables en superficie (viento y precipitación) está condicionada por el modelo conceptual que describe la evolución del ciclón. Los ciclones extratropicales intensos del Atlántico Sur suelen seguir el modelo de Fractura Frontal propuesto por \citet{shapiro1990life}. Este modelo se distingue por cuatro rasgos diagnósticos: la fractura del frente frío (\textit{frontal fracture}), el frente curvado (\textit{bent-back front}), la estructura en T (\textit{T-bone}) y el aislamiento de una masa de aire cálido en el núcleo (\textit{warm seclusion}). La observación de este modelo en casos intensos del Atlántico Sur fue documentada por \citet{shapiro1999bridge}.

% --- D2: Cinturones de transporte y organización de extremos ---
Para comprender la dinámica interna que rige la asimetría de las variables en superficie, resulta fundamental complementar este paradigma con el modelo de cinturones de transporte \citep{browning1986conceptual}. El WCB y su respectiva corriente en chorro de bajo nivel se organizan en torno al vórtice, y el ascenso y giro anticiclónico de este flujo cálido modula las estructuras de mesoescala, concentrando las bandas de precipitación y el cizallamiento del viento por delante del frente frío. Desde una perspectiva global, \citet{catto2015fronts} demuestran que hasta el 69\% de los WCBs están asociados con frentes identificables objetivamente, y que los frentes asociados a WCBs son entre 2 y 10 veces más propensos a producir precipitación extrema que los frentes sin WCB. Esta vinculación directa entre la dinámica frontal y la precipitación extrema —y no meramente la intensidad del vórtice central— justifica el enfoque espacial de esta tesis.

El acoplamiento entre WCB y CCB es especialmente relevante para explicar la asimetría conjunta de viento y precipitación. Según \citet{schemm2014linkage}, el CCB se desplaza a baja altura a lo largo del lado frío del frente curvado, experimentando un aumento de vorticidad potencial al pasar por debajo de la región de máxima liberación de calor latente generada por el WCB ascendente, localizando así el jet de bajo nivel más intenso en el extremo del frente curvado. Este acoplamiento indirecto establece que la precipitación extrema se concentra primordialmente en la zona de ascenso del WCB, mientras que los extremos de viento en superficie son modulados por la evolución del sector frío y el descenso del CCB, explicando por qué ambas variables presentan firmas espaciales asimétricas pero dinámicamente relacionadas.

% --- D3: Estructuras de mesoescala y vientos extremos ---
Dentro del paradigma de Shapiro-Keyser, una estructura de mesoescala de especial relevancia para los vientos extremos en superficie es el Sting Jet (SJ). \citet{clark2005sting} lo caracteriza como un flujo coherente que desciende desde niveles medios y produce la región de vientos superficiales más dañinos en la zona de fractura frontal del ciclón, distinta del WCB y del CCB. La distinción observacional entre el CCB y el SJ fue establecida por \citet{martinez2014cold}, quienes mediante trazadores químicos demostraron que ambas corrientes constituyen masas de aire físicamente distintas. Adicionalmente, \citet{hyeok2024extrem} identifican un mecanismo complementario: los vientos horizontales que impactan sobre la superficie frontal fría son bloqueados por esta, generando corrientes descendentes que transportan momento horizontal de niveles altos hacia la superficie, intensificando los vientos en sectores específicos del centro ciclónico. Estos mecanismos de mesoescala proveen el marco conceptual para interpretar por qué los extremos de viento en superficie pueden presentar distribuciones espaciales más complejas que las predichas por el modelo de escala sinóptica.

\subsection{Fases del ciclo de vida}\label{subsec:tb_fases}

% --- D4: Segmentación objetiva con CycloPhaser ---
Para caracterizar la evolución temporal del viento y la precipitación, este estudio adopta la segmentación del ciclo de vida en cuatro fases discretas mediante la metodología objetiva del paquete computacional \textit{CycloPhaser} \citep{deSouza2025CycloPhaser}, aplicado por \citet{coutodesouza2024new}. Este algoritmo utiliza la tasa de cambio de la vorticidad relativa $\zeta_{850}$ —conectando así directamente con el fundamento de detección establecido en la Sección~\ref{sec:tb_vorticidad}— para definir los siguientes regímenes estructurales. Cabe notar que cada fase no es homogénea internamente: los instantes próximos a sus transiciones comparten rasgos de la etapa adyacente, por lo que la condición estructural más representativa de cada régimen corresponde al núcleo temporal de la fase, alejado de ambos bordes de transición.

\begin{enumerate}
    \item Incipiente: Corresponde a la aparición inicial del núcleo de vorticidad organizada. En el contexto regional, esta fase está frecuentemente asociada a la interacción entre una vaguada en altura y forzantes de superficie, como la orografía de los Andes o gradientes térmicos costeros.

    \item Intensificación: Período caracterizado por el rápido crecimiento de la vorticidad ciclónica. Según el análisis energético de \citet{coutodesouza2024thesis}, esta etapa es impulsada por una retroalimentación entre la inestabilidad dinámica y los flujos de calor latente y sensible.

    \item Madurez: Etapa donde el sistema alcanza su intensidad máxima (pico de vorticidad) y su mayor extensión horizontal. Es durante esta fase donde los ciclones oceánicos intensos consolidan la evolución estructural documentada por \citet{shapiro1999bridge}: la deformación del flujo genera la fractura del frente frío y envuelve una masa de aire que queda aislada en el núcleo del vórtice (\textit{warm seclusion}). Al completarse este proceso, las huellas de precipitación y viento quedan espacialmente desacopladas.

    \item Decaimiento: Fase final marcada por el debilitamiento progresivo del vórtice, debido a la fricción superficial y al cese de la conversión de energía baroclínica una vez que el sistema se ha ocluido completamente o ha entrado en una zona barotrópica.
\end{enumerate}

% ============================================================
%  BLOQUE E — DISTRIBUCIÓN DE EXTREMOS EN SUPERFICIE:
%             ASIMETRÍA, EVENTOS COMPUESTOS Y HERRAMIENTAS
%             ESTADÍSTICAS
%  Propósito: Fundamentar por qué los extremos de viento y
%  precipitación son espacialmente asimétricos (E1), por qué
%  deben analizarse conjuntamente como compound events (E2),
%  y qué herramientas estadísticas son apropiadas para
%  capturar esta estructura (E3: KDE, E4: compositing,
%  E5: EOF/MEOF, E6: Bootstrap-t). Es el bloque que conecta
%  directamente con cada decisión técnica de la Metodología.
% ============================================================

\section{Estructura superficial y extremos}\label{sec:tb_compuestos}

\subsection{Asimetría espacial de los máximos}\label{subsec:tb_kde_asimetria}

% --- E1: Asimetría de precipitación y viento ---
Los ciclones extratropicales son sistemas altamente asimétricos, donde los máximos raramente coinciden con el centro de vorticidad. Es importante notar que gran parte de los modelos conceptuales clásicos fueron desarrollados para el Hemisferio Norte; por lo tanto, al analizar el Atlántico Sur, la distribución espacial debe entenderse como una imagen invertida latitudinalmente tomando al ecuador como eje.

La precipitación está fundamentalmente controlada por el WCB. En la literatura clásica del Hemisferio Norte, este flujo se describe ascendiendo típicamente hacia el noreste \citep{blanchard2021warm}; al trasladar este concepto al Hemisferio Sur, el efecto del espejo hace que el WCB concentre la humedad en el sector sureste del sistema durante su etapa de desarrollo. El ascenso inclinado de este flujo cálido organiza la precipitación en estructuras de mesoescala bien definidas, generando bandas anchas de lluvia estratiforme por delante del sistema y bandas estrechas de convección vigorosa en la interfaz del frente frío. Además, \citet{oertel2021warm} demuestran que el WCB frecuentemente alberga núcleos de convección intensa embebida. La marcada asimetría de este campo es corroborada empíricamente por \citet{naud2020evaluation}, quienes muestran que la precipitación en los ciclones oceánicos se concentra abrumadoramente en el sector oriental (ascendente y húmedo), mientras que el flanco occidental permanece subsidente y seco.

El campo de viento en superficie presenta una asimetría que cambia con el ciclo de vida del sistema. El mecanismo dinámico central que genera el jet de bajo nivel más intenso es el CCB, que —según la cadena causal descrita en la Sección~\ref{subsec:tb_modelos}— experimenta un aumento de vorticidad potencial al pasar por debajo de la región de máxima liberación de calor latente del WCB, localizando los vientos más intensos en el extremo del frente curvado. En el Atlántico Sur, \citet{bitencourt2010relating} y \citet{cardoso2022synoptic} documentan que los vientos máximos tienden a concentrarse en sectores específicos, como el cuadrante noreste, y que las ráfagas más destructivas suelen estar vinculadas a características de mesoescala cuya posición relativa al centro de vorticidad varía a medida que el ciclón madura \citep{eisenstein2023identification}. Esta heterogeneidad explica por qué la intensidad de los máximos del viento no puede capturarse mediante una única métrica central \citep{corner2025classification}.

Esta disociación entre la métrica dinámica central y el impacto en superficie tiene consecuencias directas para el diseño analítico. \citet{andrade2024composite} documentan que, en el Atlántico Sudoccidental, los ciclones más intensos según MSLP presentan trayectorias sistemáticamente desplazadas hacia el sur respecto a los más intensos según viento máximo a 10 m, y que solo el 56\% de los sistemas que superan ambos umbrales simultáneamente afectan la zona costera. Adicionalmente, \citet{gramcianinov2023impact} demuestran que los impactos más severos tienden a confinarse en sectores preferenciales del dominio centrado en el ciclón —como el oeste y noroeste en el Atlántico Sur— y que esta configuración espacial asimétrica varía fuertemente con la fase del ciclo de vida del sistema. Ambos resultados evidencian que el análisis de los extremos requiere un marco de referencia centrado en el sistema, y no un enfoque euleriano fijo.

\subsection{Eventos compuestos de viento y precipitación}\label{subsec:tb_compound_def}

% --- E2: Definición y justificación del enfoque compound ---
Históricamente, la intensidad de los ciclones se ha caracterizado mediante métricas puntuales, como la presión mínima central o la vorticidad máxima. Sin embargo, los efectos de viento y precipitación actúan frecuentemente como fenómenos meteorológicos simultáneos, y su co-ocurrencia no es estadísticamente aleatoria sino físicamente dictada por la arquitectura tridimensional del ciclón. \citet{zscheischler2018future} argumentan que los procesos que originan eventos intensos frecuentemente interactúan y presentan dependencias espaciales y temporales, lo que implica que su ocurrencia conjunta debe abordarse bajo el paradigma de \textit{compound events} (eventos compuestos).

En el contexto de esta tesis, el tipo de evento compuesto considerado es la co-ocurrencia espacial y temporal de extremos de viento y precipitación dentro de una misma fase del ciclo de vida del ciclón. Esta dependencia no es aleatoria: como se estableció en la Sección~\ref{subsec:tb_modelos}, está físicamente dictada por el acoplamiento entre el WCB (dominante en la distribución de precipitación) y el CCB (dominante en la distribución de vientos máximos). \citet{chen2025characteristics} establecen que la interacción entre el campo de viento y la precipitación no es lineal ni espacialmente homogénea, y que su ocurrencia conjunta responde a mecanismos físicos acoplados que varían según la fase de vida del sistema, lo que subraya la necesidad de analizar ambas variables en forma integrada y no por separado.

% --- E2b: Limitaciones de ERA5 para la cuantificación de extremos ---
No obstante, es necesario reconocer una limitación en la cuantificación de estas variables. \citet{chen2024evaluation} demuestran que ERA5, si bien presenta un sesgo general reducido para el viento (-0.7\%) y la precipitación (-10.4\%), este sesgo se amplifica para los eventos más intensos: para el 5\% de ciclones extratropicales más extremos, la subestimación del viento alcanza el -10.2\% y la de la precipitación el -22.6\%. Esta tendencia a suavizar los máximos más intensos implica que los resultados obtenidos para el subconjunto p90 deben interpretarse como estimaciones conservadoras de la severidad de los sistemas.

\subsection{Herramientas estadísticas para la caracterización estructural}\label{subsec:tb_herramientas}

% --- E3: Compositing centrado en el sistema ---
El análisis de la estructura típica de los ciclones requiere sintetizar información de múltiples eventos en una representación coherente. La técnica de compositing centrado en el sistema —adoptada en esta tesis sobre el marco lagrangiano descrito en la Sección~\ref{sec:tb_vorticidad}— consiste en superponer los campos de distintos eventos en coordenadas relativas al centro de cada ciclón, preservando así las asimetrías estructurales. Este enfoque opera bajo el supuesto de que los ciclones de una misma región y fase comparten un patrón estructural estadísticamente representativo, y que el promedio de sus campos lagrangianos revela dicho patrón filtrando la variabilidad idiosincrática de cada evento. Su validez metodológica ha sido demostrada por estudios como el de \citet{priestley2022improved}, quienes aplican compositing centrado al 10\% de ciclones más intensos para aislar estructuras representativas, y por \citet{han2025system}, quienes muestran que esta estrategia es capaz de revelar cómo los campos extremos cambian de ubicación o extensión a lo largo del ciclo de vida del sistema.

% --- E4: Estimación de densidad por kernel (KDE) ---
Para caracterizar la distribución espacial preferencial de los máximos de cada variable sin imponer supuestos distribucionales a priori, se adopta la Estimación de Densidad por Kernel (KDE). Esta técnica transforma observaciones discretas de posición —la ubicación relativa del máximo de cada ciclón respecto a su centro— en una superficie de probabilidad continua, permitiendo mapear la huella espacial del evento sin necesidad de definir celdas o grillas de conteo. Como señala \citet{weglarczyk2018kernel}, la estimación por núcleos elimina el sesgo inherente a la selección del número, ancho y punto de inicio de las clases de un histograma, ya que preserva y utiliza la ubicación exacta de cada observación. En su extensión bidimensional, la técnica genera un campo continuo de densidad $\hat{f}(\mathbf{s})$ que cuantifica, para cada posición $\mathbf{s} = (\Delta x, \Delta y)$ dentro del dominio centrado, la tendencia espacial a encontrar el máximo de la variable. La formulación matemática detallada de esta estimación y los parámetros de implementación específicos se desarrollan en la Sección~\ref{subsec:kde_localizacion} de la Metodología.

% --- E5: Descomposición modal EOF y MEOF ---
Para aislar los patrones espaciales coherentes de variabilidad dentro del acoplamiento físico entre viento y precipitación, el análisis de Funciones Ortogonales Empíricas (EOF) y su extensión multivariada (MEOF) se han establecido como el estándar metodológico. Según \citet{hannachi2007empirical}, esta técnica permite descomponer la variabilidad de un campo en modos ortogonales de varianza máxima, separando la señal física robusta del ruido de mesoescala. Específicamente para eventos compuestos, \citet{liang2018multivariate} demostraron la eficacia del MEOF al aplicarlo sobre un vector de estado normalizado que integra métricas con unidades físicas dispares (como la precipitación y el viento superficial). Esta normalización estadística permite identificar si los máximos de ambas variables covarían en el mismo espacio geométrico o si, por el contrario, presentan desfases estructurales sistemáticos. La formulación matemática detallada, el procedimiento de normalización y la implementación computacional específica para este estudio se desarrollan exhaustivamente en la Sección~\ref{subsec:eof_multivariado} de la Metodología.

% --- E6: Significancia estadística mediante Bootstrap-t ---
Para evaluar la robustez estadística de los patrones espaciales derivados del análisis EOF y aislar las señales físicas del ruido estocástico, se adopta el método de remuestreo no paramétrico Bootstrap-t \citep{efron1993introduction}. A diferencia de las pruebas paramétricas estándar —que asumen distribuciones normales en las variables—, el Bootstrap-t no impone supuestos distribucionales, lo que lo hace especialmente apropiado para las distribuciones de precipitación y viento ciclónico, que son intrínsecamente asimétricas y de cola pesada. Como establece \citet{hesterberg2015bootstrap}, la variante Bootstrap-t corrige dinámicamente esta asimetría al estandarizar cada muestra de remuestreo por su propio error estándar, generando estadísticos cuasi-pivotales que producen intervalos de confianza con cobertura más precisa que el Bootstrap percentílico estándar. La implementación específica de este procedimiento, incluyendo el número de iteraciones y los criterios de significancia aplicados a cada punto de grilla, se detalla en la Sección~\ref{subsec:bootstrap_significancia} de la Metodología.